%% kapa: exam=1 round=22 year=2011 period=2 form=A subject=accounting subject_ko=회계학 seq=38 number=38
%% source: Q-Net 감정평가사 기출문제 (cst00309) artlSeq=5036790 file=extracted/1008/5036790/2011년 제22회 감정평가사 1차시험 기출문제/2교시.hwp
%% conversion: kordoc 4.24.51 -> tools/gen1.py (auto); review=proofread-sonnet; flags=none
\begin{question}{38}
(주)서울은 기계A나 기계B를 구입하여 신형자전거를 생산하려고 한다. 신형자전거가 생산되면 구매자의 선호에 따라 히트상품이 될 수도 있고 보통상품이 될 수도 있다. 각 상황에 따라 예상되는 이익은 다음과 같다.

{\small\centering
\begin{tabular}{|l|c|c|}
\hline
\multirow{2}{*}{구 분} & \multicolumn{2}{c|}{미래상황} \\
\cline{2-3}
 & 히트상품 & 보통상품 \\
\hline
기계A 구입 & ￦120,000 & ￦40,000 \\
기계B 구입 & 150,000 & 10,000 \\
\hline
\end{tabular}\par}

신형자전거가 히트상품이 될 확률이 40\%이며 보통상품이 될 확률은 60\%라고 한다. 다음 중 옳지 않은 것은? \par
\begin{choices}
  \item 기계A를 구입하는 대안의 기대이익은 ￦72,000이다.
  \item 기계B를 구입하고 신형자전거가 보통상품이 될 경우 조건부 손실(conditional loss)은 ￦30,000이다.
  \item 각 상황에 대해 80\% 정확도를 가진 보고서가 있다면, 이 보고서는 정보로서의 가치가 있다.
  \item 각 상황에 대해 100\% 정확한 예측을 하는 보고서가 있을 때, 이 보고서의 최대 구입가격은 ￦12,000이다.
  \item 조건부 손실의 기대값을 최소화하는 대안은 기계B를 구입하는 것이다.
\end{choices}
\end{question}
\answer{5}
